\chapter{Worth Revisiting}
\label{ch:revisit}

Thirteen things from the week that reward a second pass, in no particular order. Some are readings,
some are exercises, and a few are loose ends that nobody has tied up.

\begin{enumerate}

\item \textbf{The 4-bit failure in the quantization tutorial}, in the section headed 8-bit RTN.
  Round-to-nearest quantization is essentially lossless at 8 bits and visibly destroys the same model
  at 4. Watching that happen before reading Chapter~\ref{ch:10-efficient-learning} is what makes the
  inverse-Hessian derivation feel necessary rather than decorative.

\item \textbf{Activation steering}, exercises 9 to 11 of the mechanistic interpretability tutorial.
  A steering vector is built from the contrast between two sets of prompts, added to the residual
  stream mid-generation, and the effect on the output is unmistakable. The whole method fits in three
  exercises and transfers to any model that can be wrapped in TransformerLens.

\item \textbf{Continual learning, read as a disagreement.} Chapter~\ref{ch:06-continual-learning}
  and then Chapter~\ref{ch:07-to-be-figured-out}. Lyle treats plasticity loss as a fixable fault in
  the optimiser and the architecture; Pascanu treats the same territory as evidence that
  gradient-based credit assignment structurally depends on i.i.d.\ data. Reading them in that order
  makes the gap between the two positions visible in a way neither lecture does alone.

\item \textbf{Dieleman on diffusion as spectral autoregression},
  \url{https://sander.ai/2023/07/20/perspectives.html}. The slide version, Figure~\ref{fig:diffusion-models-3}, compresses
  the argument to four panels. The post carries the derivation, and it is the most reusable idea in
  the diffusion lecture: the noise schedule is a frequency schedule.

\item \textbf{PPO and GRPO, implemented in that order}, Part~2 of the reinforcement learning
  tutorial. PPO is built first with a critic, at full memory cost, and GRPO then removes the critic
  in a few lines by standardising rewards within a sampled group. The ordering is what makes GRPO's
  appeal legible rather than merely stated.

\item \textbf{The Kazhdan--Lusztig experiment}~\citep{davies2021}, with
  Chapter~\ref{ch:08-ml-for-mathematics}. Gradient attribution on a trained GNN flagged particular
  edges, those edges suggested a decomposition, and the decomposition became a conjecture. It is the
  clearest example available of a model's internals being legible enough to turn into a mathematical
  statement.

\item \textbf{The linear recurrent unit derivation, all six steps.} \citet{orvieto2023}, with
  Chapter~\ref{ch:07-to-be-figured-out}. An ablation dressed as an architecture, where the negative
  results are the substance: HiPPO is unnecessary, continuous time is unnecessary, and complex
  numbers may buy nothing beyond compactness. A better first reading on state space models than most
  explainers.

\item \textbf{CKA and the invariance argument around it}~\citep{kornblith2019}, with
  Chapter~\ref{ch:02-representational-geometry}. A small formula with a great deal of thinking behind
  it. The transferable habit is to ask of any similarity measure which transformations it has been
  built to ignore, and whether those are the ones the downstream task ignores too.

\item \textbf{The eight-episode TD versus Monte Carlo exercise.} Slides 37 to 39 of the
  reinforcement learning deck, with Chapter~\ref{ch:04-intro-reinforcement-learning}. Five minutes
  with pen and paper, and it fixes the difference between fitting the observed data and fitting the
  model that data implies. Best attempted before looking at the answer.

\item \textbf{Computational identifiability.} Chapter~\ref{ch:05-statistics-estimation}, final
  section, and~\citet{bynum2026identifiability}. Identifiability defined relative to a prior over
  problems, a hypothesis space and an optimiser, rather than absolutely. The idea from the week most
  likely to change how a problem is framed rather than how it is solved.

\item \textbf{Zero-padding applied to causal attention.}
  Chapter~\ref{ch:09-spectral-analysis-dags}, with the inferred direction in
  Chapter~\ref{ch:research-directions}. A causally-masked transformer is a DAG, so its adjacency
  spectrum collapses. Whether Stankovi\'{c}'s construction says anything useful about attention is
  untested, and cheap to sanity-check on a small mask.

\item \textbf{PECNET with a state space model}, presented at the poster session by Mustafa Abdullah
  Hakkoz for regional agricultural forecasting. Worth pairing with the state space material in
  Chapter~\ref{ch:07-to-be-figured-out}, since the poster puts a hybrid to work on a real forecasting
  problem rather than on a benchmark.

\item \textbf{The geometric deep learning prerequisites},
  \url{https://arxiv.org/abs/2508.02723}, with the course at
  \url{https://geometricdeeplearning.com}. The group- and representation-theoretic background is not
  needed to apply the four Gs, only to derive them. Time well spent for anyone intending to build an
  architecture from a symmetry rather than pick one off the shelf.

\end{enumerate}
